\begin{table}[htbp]\centering\small
\caption{Ingresos del artículo 1o., línea P: proyecto 2026 contra proyecto 2027 (mdp)}\label{cua:ing}
\begin{tabular}{lrrrr}
\toprule
Concepto & 2026 & 2027 & Var. nominal (\%) & Var. real (\%) \\
\midrule
Total & 10.193.683,7 & 10.636.488,1 & +4,3 & +0,3 \\
Impuestos & 5.838.541,1 & 6.263.886,8 & +7,3 & +3,2 \\
\quad Impuesto sobre la renta & 3.070.149,1 & 3.288.975,3 & +7,1 & +3,0 \\
\quad Impuesto al valor agregado & 1.589.069,0 & 1.768.017,3 & +11,3 & +7,0 \\
\quad IEPS & 761.501,9 & 828.256,0 & +8,8 & +4,6 \\
\quad\quad de combustibles automotrices & 473.279,1 & 538.549,2 & +13,8 & +9,4 \\
Cuotas de seguridad social & 641.782,1 & 702.564,7 & +9,5 & +5,3 \\
Derechos & 157.081,7 & 177.185,9 & +12,8 & +8,5 \\
Aprovechamientos & 203.520,5 & 329.479,8 & +61,9 & +55,7 \\
Ventas de bienes y servicios & 1.630.973,6 & 1.457.539,9 & $-$10,6 & $-$14,1 \\
Transferencias del FMP & 232.630,4 & 210.774,0 & $-$9,4 & $-$12,9 \\
Ingresos derivados de financiamientos & 1.472.626,4 & 1.479.959,2 & +0,5 & $-$3,4 \\
\bottomrule
\end{tabular}
\\[2pt]\begin{minipage}{0.92\textwidth}\scriptsize Fuente: artículo 1o. de la ILIF 2027 y de la ILIF 2026. Deflactor del PIB de 2027, 1,040. \textbf{El renglón del IEPS de 2026 se restituye de la LIF aprobada}, porque el diseño de la tabla de la iniciativa no lo aísla; iniciativa y ley aprobada de 2026 coinciden al décimo en todos los demás renglones de este cuadro, incluido el total, de modo que para 2026 las líneas P y G son la misma.\end{minipage}
\end{table}
